% Appendix shared by the competitor manuals: the starter package, file by file, and the
% driver's logic per scan. Needs tikz and listings loaded by the including document.

\section{The starter package, file by file}
\label{app:starter}

\texttt{starter\_kit/feb\_driver} is a normal ROS~2 Python package. After
\texttt{cp -r starter\_kit/feb\_driver stack/} these are the files you own (the Python files live in the inner \texttt{feb\_driver/} folder):

\begin{tabularx}{\linewidth}{p{4.4cm}X}
\toprule
File & What it is \\
\midrule
\texttt{package.xml} & The ROS~2 manifest: the package name \texttt{feb\_driver}, its version and
its dependencies (\texttt{rclpy}, \texttt{sensor\_msgs}, \texttt{std\_msgs}). \texttt{colcon}
reads it to know the package exists. You only touch it to add a dependency. \\
\texttt{setup.py} & Python packaging. The line that matters is the entry point
\texttt{reactive\_driver = feb\_driver.reactive\_driver:main}: it is what makes
\texttt{ros2 run feb\_driver reactive\_driver} exist. It also installs \texttt{launch/} and
\texttt{config/} so launch files can find them. Add a line here for every new node. \\
\texttt{setup.cfg} & Two lines telling \texttt{colcon} where to put the installed scripts. Never edited. \\
\texttt{resource/feb\_driver} & An empty marker file ROS~2 uses to index the package. Never edited. \\
\texttt{\_\_init\_\_.py} & Empty; makes the folder a Python module. \\
\texttt{reactive\_driver.py} & The driver, about 170 lines, explained below. \\
\texttt{config/driver.yaml} & Every tuning number, one per line with a comment. The launch file
hands it to the node, so you change behaviour here without touching code. \\
\texttt{launch/drive.launch.py} & Starts one node, \texttt{reactive\_driver}, with the yaml and
\texttt{use\_sim\_time} set, so the node's clock follows the simulator. This is what
\texttt{--stack "ros2 launch feb\_driver drive.launch.py"} runs. \\
\texttt{sysid\_probe.py}, \texttt{launch/sysid.launch.py} & Stage 2. Drives with the
same steering but a fixed throttle schedule (three steps, a coast, a chirp) while recording a
bag, to measure how the car responds. \\
\texttt{tools/sysid\_fit.py} & Fits that bag: steady-state gain (the \texttt{speed\_per\_throttle}
number), idle brake, throttle lag and effective wheelbase. Run inside the container. \\
\bottomrule
\end{tabularx}

\medskip
The node subscribes to three topics and publishes two, all under
\texttt{/autodrive/roboracer\_1/}: in, \texttt{lidar} (\texttt{LaserScan}, 1081 ranges over
$\pm135^\circ$) and \texttt{left\_encoder}, \texttt{right\_encoder} (\texttt{JointState}, wheel
angle); out, \texttt{steering\_command} and \texttt{throttle\_command} (\texttt{Float32}).
Steering is normalised, $-1$ to $1$ meaning $-30^\circ$ to $30^\circ$. Throttle is $0$ to $1$
and \textbf{0 is a hard brake}, not coasting.

\section{How the driver drives, one scan at a time}
\label{app:driver}

Every lidar scan, about ten times a second, \texttt{on\_scan} runs the steps below and
publishes one steering and one throttle value. The figure shows one scan from above, in a corridor that bends left; the right wall is the nearest thing the lidar sees.

\begin{center}
\begin{tikzpicture}[scale=1.1, every node/.style={font=\small}]
  \begin{scope}
  \clip (-4.0,-1.5) rectangle (4.4,5.0);
  \draw[line width=1.8pt, gray!75] (1.000,-2.800) -- (1.000,2.200) -- (1.000,2.200) -- (0.994,2.370) -- (0.978,2.539) -- (0.950,2.707) -- (0.911,2.873) -- (0.862,3.036) -- (0.802,3.195) -- (0.732,3.350) -- (0.652,3.500) -- (0.562,3.644) -- (0.463,3.783) -- (0.355,3.914) -- (0.238,4.038) -- (0.114,4.155) -- (-0.017,4.263) -- (-0.156,4.362) -- (-0.300,4.452) -- (-0.450,4.532) -- (-0.605,4.602) -- (-0.764,4.662) -- (-0.927,4.711) -- (-1.093,4.750) -- (-1.261,4.778) -- (-1.430,4.794) -- (-1.600,4.800);
  \draw[line width=1.8pt, gray!75] (-1.600,-2.800) -- (-1.600,1.400) -- (-1.600,1.400) -- (-1.602,1.452) -- (-1.607,1.504) -- (-1.615,1.556) -- (-1.627,1.607) -- (-1.642,1.657) -- (-1.661,1.706) -- (-1.683,1.754) -- (-1.707,1.800) -- (-1.735,1.844) -- (-1.765,1.887) -- (-1.799,1.927) -- (-1.834,1.966) -- (-1.873,2.001) -- (-1.913,2.035) -- (-1.956,2.065) -- (-2.000,2.093) -- (-2.046,2.117) -- (-2.094,2.139) -- (-2.143,2.158) -- (-2.193,2.173) -- (-2.244,2.185) -- (-2.296,2.193) -- (-2.348,2.198) -- (-2.400,2.200);
  \fill[febblue, opacity=0.13] (0,0) -- (-1.600,0.000) -- (-1.600,0.140) -- (-1.600,0.282) -- (-1.600,0.429) -- (-1.600,0.582) -- (-1.600,0.746) -- (-1.600,0.924) -- (-1.600,1.120) -- (-1.600,1.343) -- (-1.636,1.636) -- (-3.857,4.596) -- (-3.441,4.915) -- (-3.000,5.196) -- (-2.536,5.438) -- (-2.052,5.638) -- (-1.281,4.780) -- (-0.825,4.681) -- (-0.394,4.502) -- (0.000,4.249) -- (0.344,3.926) -- (0.625,3.543) -- (0.834,3.111) -- (0.961,2.640) -- (1.000,2.145) -- (1.000,1.732) -- (1.000,1.428) -- (1.000,1.192) -- (1.000,1.000) -- (1.000,0.839) -- (1.000,0.700) -- (1.000,0.577) -- (1.000,0.466) -- cycle;
  \draw[gray!30, very thin] (0,0) -- (-1.600,-1.600);
  \draw[gray!30, very thin] (0,0) -- (-1.600,-1.343);
  \draw[gray!30, very thin] (0,0) -- (-1.600,-1.120);
  \draw[gray!30, very thin] (0,0) -- (-1.600,-0.924);
  \draw[gray!30, very thin] (0,0) -- (-1.600,-0.746);
  \draw[gray!30, very thin] (0,0) -- (-1.600,-0.582);
  \draw[gray!30, very thin] (0,0) -- (-1.600,-0.429);
  \draw[gray!30, very thin] (0,0) -- (-1.600,-0.282);
  \draw[gray!30, very thin] (0,0) -- (-1.600,-0.140);
  \draw[gray!55, thin] (0,0) -- (-1.600,0.000);
  \draw[gray!55, thin] (0,0) -- (-1.600,0.140);
  \draw[gray!55, thin] (0,0) -- (-1.600,0.282);
  \draw[gray!55, thin] (0,0) -- (-1.600,0.429);
  \draw[gray!55, thin] (0,0) -- (-1.600,0.582);
  \draw[gray!55, thin] (0,0) -- (-1.600,0.746);
  \draw[gray!55, thin] (0,0) -- (-1.600,0.924);
  \draw[gray!55, thin] (0,0) -- (-1.600,1.120);
  \draw[gray!55, thin] (0,0) -- (-1.600,1.343);
  \draw[gray!55, thin] (0,0) -- (-1.636,1.636);
  \draw[gray!55, thin] (0,0) -- (-3.857,4.596);
  \draw[gray!55, thin] (0,0) -- (-3.441,4.915);
  \draw[gray!55, thin] (0,0) -- (-3.000,5.196);
  \draw[gray!55, thin] (0,0) -- (-2.536,5.438);
  \draw[gray!55, thin] (0,0) -- (-2.052,5.638);
  \draw[gray!55, thin] (0,0) -- (-1.281,4.780);
  \draw[gray!55, thin] (0,0) -- (-0.825,4.681);
  \draw[gray!55, thin] (0,0) -- (-0.394,4.502);
  \draw[gray!55, thin] (0,0) -- (0.000,4.249);
  \draw[gray!55, thin] (0,0) -- (0.344,3.926);
  \draw[gray!55, thin] (0,0) -- (0.625,3.543);
  \draw[gray!55, thin] (0,0) -- (0.834,3.111);
  \draw[gray!55, thin] (0,0) -- (0.961,2.640);
  \draw[gray!55, thin] (0,0) -- (1.000,2.145);
  \draw[gray!55, thin] (0,0) -- (1.000,1.732);
  \draw[gray!55, thin] (0,0) -- (1.000,1.428);
  \draw[gray!55, thin] (0,0) -- (1.000,1.192);
  \draw[gray!55, thin] (0,0) -- (1.000,1.000);
  \draw[gray!55, thin] (0,0) -- (1.000,0.839);
  \draw[gray!55, thin] (0,0) -- (1.000,0.700);
  \draw[gray!55, thin] (0,0) -- (1.000,0.577);
  \draw[gray!55, thin] (0,0) -- (1.000,0.466);
  \draw[gray!30, very thin] (0,0) -- (1.000,-0.087);
  \draw[gray!30, very thin] (0,0) -- (1.000,-0.176);
  \draw[gray!30, very thin] (0,0) -- (1.000,-0.268);
  \draw[gray!30, very thin] (0,0) -- (1.000,-0.364);
  \draw[gray!30, very thin] (0,0) -- (1.000,-0.466);
  \draw[gray!30, very thin] (0,0) -- (1.000,-0.577);
  \draw[gray!30, very thin] (0,0) -- (1.000,-0.700);
  \draw[gray!30, very thin] (0,0) -- (1.000,-0.839);
  \draw[gray!30, very thin] (0,0) -- (1.000,-1.000);
  \draw[red!70, thick] (0,0) -- (1.000,0.364);
  \draw[red!70, thick] (0,0) -- (1.000,0.268);
  \draw[red!70, thick] (0,0) -- (1.000,0.176);
  \draw[red!70, thick] (0,0) -- (1.000,0.087);
  \draw[red!70, thick] (0,0) -- (1.000,0.000);
  \fill[red] (1.000,0.000) circle (2.4pt);
  \draw[red, dashed, thick] (1.000,0.000) circle (0.400);
  \draw[febblue, line width=1.6pt] (0,0) -- (-3.857,4.596);
  \draw[febblue!80, thick, dash pattern=on 3pt off 2pt] (0,0) -- (-0.825,4.681);
  \draw[-{Stealth[length=8pt]}, line width=2pt, black] (0,0) -- (-1.080,2.031);
  \draw[green!45!black, dashed] (-0.235,0) -- (-0.235,4.15); \draw[green!45!black, dashed] (0.235,0) -- (0.235,4.15);
  \fill[black!80] (-0.135,-0.27) rectangle (0.135,0.27);
  \end{scope}
  \node[anchor=west] (n) at (2.90,-0.20) {nearest return, 1.0 m};
  \draw[gray!60, thin] (n.west) -- (1.000,0.000);
  \node[anchor=west] (n) at (2.90,-0.60) {bubble 0.4 m: beams within $\pm22^\circ$ blanked};
  \draw[gray!60, thin] (n.west) -- (0.700,0.000);
  \node[anchor=east] (n) at (-3.60,1.00) {widest gap};
  \draw[gray!60, thin] (n.east) -- (-0.960,0.554);
  \node[anchor=east] (n) at (-3.60,3.90) {deepest beam};
  \draw[gray!60, thin] (n.east) -- (-3.278,3.907);
  \node[anchor=east] (n) at (-3.60,3.30) {gap middle};
  \draw[gray!60, thin] (n.east) -- (-0.578,3.276);
  \node[anchor=west] (n) at (2.90,2.90) {target: 0.6 deepest + 0.4 middle};
  \draw[gray!60, thin] (n.west) -- (-1.080,2.031);
  \node[anchor=west] (n) at (2.90,2.20) {room ahead: 3rd percentile of ranges};
  \draw[gray!60, thin] (n.west) -- (0.235, 3.748586362432566);
  \node[anchor=west] (n) at (2.90,1.90) {inside the car's own 0.47 m corridor};
  \node[anchor=west] (n) at (2.90,0.40) {car, lidar at the origin};
  \draw[gray!60, thin] (n.west) -- (0.135, 0.0);
  \node[anchor=west] (n) at (2.90,-1.20) {beams behind the car: outside the gap search};
  \draw[gray!60, thin] (n.west) -- (0.700,-0.404);
\end{tikzpicture}
\end{center}

\begin{enumerate}[leftmargin=*]
  \item \textbf{Clean the scan.} Any beam that returned nothing (infinity) is set to the lidar's
  maximum range, and every range is clipped to the sensor's limits. This is the numpy array
  \texttt{ranges}; \texttt{angles} is the matching bearing of each beam.
  \item \textbf{Find the nearest return and blank a bubble around it.} Only beams within
  $\pm$\texttt{gap\_fov} (1.6 rad, about $\pm 92^\circ$) count. The shortest of them is the
  nearest obstacle, the red dot. Every beam whose bearing is within
  $\arctan(\texttt{bubble}/r_{\min})$ of that beam is set to zero: at $r_{\min}=1.0$~m and
  \texttt{bubble}$=0.4$~m that is $\pm 22^\circ$, the red beams. A zeroed beam can never be
  chosen as a direction, so the car never aims through a gap that passes within 0.4~m of the
  closest thing it sees. That is the whole safety idea, and it is why the car brushes walls when
  \texttt{bubble} is too small and refuses corners when it is too big.
  \item \textbf{Pick the widest gap.} The non-zero beams are split into runs of consecutive
  beams; the longest run is the gap (blue). Inside it, two candidates: the deepest beam (the
  longest range) and the middle beam of the run. The target bearing is a blend,
  \texttt{deep\_weight} $\times$ deepest $+$ $(1-\texttt{deep\_weight})$ $\times$ middle, $0.6/0.4$
  by default. Pure deepest cuts corners tight and hunts; pure middle is smooth but wide.
  \item \textbf{Turn the bearing into steering.} The car is still turning at its current
  steering angle for the actuator delay (0.25~s) plus one scan period before the new command
  takes effect, so the heading change over that time is subtracted from the target first
  (delay compensation). Then \texttt{steer\_gain} scales the bearing (0.7: a $20^\circ$ target
  asks for $14^\circ$ of steering), the result is clipped to the $\pm 30^\circ$ limit, low-passed
  with \texttt{steer\_tau} and rate-limited to the actuator's 3.2~rad/s. Without the filtering,
  a target that jumps from one beam to the next would flick the wheel full lock every scan.
  \item \textbf{Measure the room ahead.} Take the beams whose end point lies inside the car's
  own corridor (a strip 0.235~m either side of the centreline, the green dashes). The 3rd
  percentile of their forward distance is the clearance: not the minimum, which one noisy beam
  would spoil, but close to it. It is low-passed, except that a genuinely close obstacle
  (under \texttt{brake\_margin} $+0.3$~m) passes straight through so braking is immediate.
  \item \textbf{Pick a speed.} Three ceilings: \texttt{max\_speed} (2.5~m/s); the cornering
  speed $\sqrt{\texttt{lat\_accel}/\kappa}$ from the slowly filtered steering, where
  $\kappa=\tan(\delta)/L$ is the path curvature; and the braking speed
  $\sqrt{2\,\texttt{decel\_limit}\,(\mathrm{clearance}-\texttt{brake\_margin})}$, the fastest speed
  from which the car can still stop before the clearance. The goal is the smallest of the
  three, floored at \texttt{min\_speed} (0.9~m/s) so the car never parks in a corner. The goal
  is low-passed and then slewed at \texttt{accel\_limit} up and \texttt{decel\_limit} down.
  \item \textbf{Speed to throttle.} The wheel encoders give the measured speed (wheel angle
  change over the last 0.25~s times the wheel radius). Throttle is feed-forward plus trim:
  feed-forward is $v_{\mathrm{target}}/\texttt{speed\_per\_throttle}$ (23 m/s per unit throttle,
  measured with the sysid tools, so 1~m/s needs about 0.043); trim is a small PI on the speed
  error, clipped to half of the feed-forward. Because of that clip the command can never reach
  zero while speed is wanted, and zero would be a hard brake. Finally the throttle is
  slew-limited to 0.8 per second. A textbook PI speed loop without these two guards surges and
  brakes in a 0.8~s cycle on this simulator; this one holds speed at 10~Hz.
  \item \textbf{Publish.} Steering divided by $30^\circ$ goes to \texttt{steering\_command},
  the throttle to \texttt{throttle\_command}. On shutdown the node publishes $0, 0$ once so the
  car stops.
\end{enumerate}

\textbf{Follow mode} (\texttt{follow\_gap} $> 0$, used by the house racer as the opponent):
if something sits within that distance in a $\pm 15^\circ$ cone straight ahead, it is treated
as a car to sit behind, not a wall to steer around. Its beams are replaced by the corridor seen
beside it, so the steering keeps the corridor's line, and the speed goal is capped to hold the
distance.

\textbf{Where to start tuning.} \texttt{bubble} and \texttt{deep\_weight} decide the line;
\texttt{max\_speed}, \texttt{lat\_accel} and \texttt{brake\_margin} decide how fast it dares.
Change one number, run the same command, watch the lap time. The sysid tools are for when
you want the throttle law to match the car rather than guess.
